\begin{sagesilent}
def extract_parameters_Y_s_p(V):
    t = var('t')
    resultats_Y = []
    resultats_s_p = []

    # On vérifie si V est une somme.
    # Si operator(V) est l'addition, on prend les opérandes, sinon on met V dans une liste.
    if hasattr(V, 'operator') and V.operator() is not None and "add" in str(V.operator()):
        termes = V.operands()
    else:
        termes = [V]

    for terme in termes:
        partie_exp = None

        # 1. On cherche si le terme contient une exponentielle
        # On regarde le terme lui-même et ses sous-parties
        if "exp" in str(terme.operator()):
            partie_exp = terme
        else:
            for op in terme.operands():
                if "exp" in str(op.operator()):
                    partie_exp = op
                    break

        # 2. Extraction des valeurs
        if partie_exp is not None:
            # s_p est l'argument de l'exp (1er opérande) divisé par t
            s_p = partie_exp.operands()[0] / t
            # Y est le terme total divisé par l'exponentielle
            Y = terme / partie_exp

            resultats_Y.append(Y)
            resultats_s_p.append(s_p)
        else:
            # Si aucune exp n'est trouvée (cas du 6 par exemple)
            resultats_Y.append(terme)
            resultats_s_p.append(0)

    return resultats_Y, resultats_s_p
\end{sagesilent}
